\begin{thebibliography}{00}

\bibitem{imran2020}
A. Imran et al., ``AI4COVID-19: AI enabled preliminary diagnosis for COVID-19 from cough samples via an app,'' \emph{Inform. Med. Unlocked}, vol. 20, p. 100378, 2020.

\bibitem{laguarta2020}
J. Laguarta, F. Hueto, and B. Subirana, ``COVID-19 artificial intelligence diagnosis using only cough recordings,'' \emph{IEEE Open J. Eng. Med. Biol.}, vol. 1, pp. 275--281, 2020.

\bibitem{pahar2021}
M. Pahar, M. Klopper, R. Warren, and T. Niesler, ``COVID-19 cough classification using machine learning and global smartphone recordings,'' \emph{Comput. Biol. Med.}, vol. 135, p. 104572, 2021.

\bibitem{roberts2021}
M. Roberts et al., ``Common pitfalls and recommendations for using machine learning to detect and prognosticate for COVID-19 using chest radiographs and CT scans,'' \emph{Nature Mach. Intell.}, vol. 3, no. 3, pp. 199--217, 2021.

\bibitem{han2022}
J. Han et al., ``Sounds of COVID-19: Exploring realistic performance of audio-based digital testing,'' \emph{npj Digit. Med.}, vol. 5, p. 16, 2022.

\bibitem{coppock2024}
H. Coppock et al., ``Audio-based AI classifiers show no evidence of improved COVID-19 screening over simple symptoms checkers,'' \emph{Nature Mach. Intell.}, vol. 6, pp. 229--242, 2024.

\bibitem{xia2021}
T. Xia et al., ``COVID-19 sounds: A large-scale audio dataset for digital respiratory screening,'' in \emph{Adv. Neural Inf. Process. Syst. 34, Datasets and Benchmarks Track}, 2021.

\bibitem{zhou2021}
Q. Zhou et al., ``Cough recognition based on mel-spectrogram and convolutional neural network,'' \emph{Frontiers in Robotics and AI}, vol. 8, p. 580080, 2021.

\bibitem{coswara2020}
N. Sharma et al., ``Coswara: A database of breathing, cough, and voice sounds for COVID-19 diagnosis,'' in \emph{Proc. Interspeech}, 2020, pp. 4811--4815.

\bibitem{fsd50k2022}
E. Fonseca, X. Favory, J. Pons, F. Font, and X. Serra, ``FSD50K: An open dataset of human-labeled sound events,'' \emph{IEEE/ACM Trans. Audio, Speech, Lang. Process.}, vol. 30, pp. 829--852, 2022.

\bibitem{coughvid2021}
L. Orlandic, T. Teijeiro, and D. Atienza, ``The COUGHVID crowdsourcing dataset, a corpus for the study of large-scale cough analysis algorithms,'' \emph{Scientific Data}, vol. 8, p. 156, 2021.

\bibitem{gemmeke2017}
J. F. Gemmeke et al., ``Audio Set: An ontology and human-labeled dataset for audio events,'' in \emph{Proc. IEEE Int. Conf. Acoust., Speech Signal Process. (ICASSP)}, 2017, pp. 776--780.

\bibitem{he2016}
K. He, X. Zhang, S. Ren, and J. Sun, ``Deep residual learning for image recognition,'' in \emph{Proc. IEEE Conf. Comput. Vis. Pattern Recognit. (CVPR)}, 2016, pp. 770--778.

\bibitem{sandler2018}
M. Sandler, A. Howard, M. Zhu, A. Zhmoginov, and L.-C. Chen, ``MobileNetV2: Inverted residuals and linear bottlenecks,'' in \emph{Proc. IEEE Conf. Comput. Vis. Pattern Recognit. (CVPR)}, 2018, pp. 4510--4520.

\bibitem{huang2017}
G. Huang, Z. Liu, L. van der Maaten, and K. Q. Weinberger, ``Densely connected convolutional networks,'' in \emph{Proc. IEEE Conf. Comput. Vis. Pattern Recognit. (CVPR)}, 2017, pp. 2261--2269.

\bibitem{zhang2026}
W. Zhang, T. Teijeiro, J. Thevenot, and D. Atienza, ``Why ML-based cough models do not generalize: A systematic cross-dataset evaluation for tuberculosis screening,'' 2026, arXiv:2608.25846.

\bibitem{vocalsound2022}
Y. Gong, J. Yu, and J. Glass, ``VocalSound: A dataset for improving human vocal sounds recognition,'' in \emph{Proc. IEEE Int. Conf. Acoust., Speech Signal Process. (ICASSP)}, 2022.

\bibitem{donoho1994}
D. L. Donoho and I. M. Johnstone, ``Ideal spatial adaptation by wavelet shrinkage,'' \emph{Biometrika}, vol. 81, no. 3, pp. 425--455, 1994.

\bibitem{chang2000}
S. G. Chang, B. Yu, and M. Vetterli, ``Adaptive wavelet thresholding for image denoising and compression,'' \emph{IEEE Trans. Image Process.}, vol. 9, no. 9, pp. 1532--1546, 2000.

\bibitem{johnstone2005}
I. M. Johnstone and B. W. Silverman, ``Empirical Bayes selection of wavelet thresholds,'' \emph{Ann. Statist.}, vol. 33, no. 4, pp. 1700--1752, 2005.

\bibitem{mallat1989}
S. G. Mallat, ``A theory for multiresolution signal decomposition: The wavelet representation,'' \emph{IEEE Trans. Pattern Anal. Mach. Intell.}, vol. 11, no. 7, pp. 674--693, 1989.

\bibitem{stevens1937}
S. S. Stevens, J. Volkmann, and E. B. Newman, ``A scale for the measurement of the psychological magnitude pitch,'' \emph{J. Acoust. Soc. Amer.}, vol. 8, no. 3, pp. 185--190, 1937.

\bibitem{russakovsky2015}
O. Russakovsky et al., ``ImageNet large scale visual recognition challenge,'' \emph{Int. J. Comput. Vis.}, vol. 115, no. 3, pp. 211--252, 2015.

\bibitem{kingma2015}
D. P. Kingma and J. Ba, ``Adam: A method for stochastic optimization,'' in \emph{Proc. 3rd Int. Conf. Learn. Representations (ICLR)}, 2015.

\bibitem{paszke2019}
A. Paszke et al., ``PyTorch: An imperative style, high-performance deep learning library,'' in \emph{Adv. Neural Inf. Process. Syst. 32 (NeurIPS)}, 2019.

\bibitem{matthews1975}
B. W. Matthews, ``Comparison of the predicted and observed secondary structure of T4 phage lysozyme,'' \emph{Biochimica et Biophysica Acta}, vol. 405, no. 2, pp. 442--451, 1975.

\bibitem{chicco2020}
D. Chicco and G. Jurman, ``The advantages of the Matthews correlation coefficient (MCC) over F1 score and accuracy in binary classification evaluation,'' \emph{BMC Genomics}, vol. 21, p. 6, 2020.

\bibitem{wilson1927}
E. B. Wilson, ``Probable inference, the law of succession, and statistical inference,'' \emph{J. Amer. Statist. Assoc.}, vol. 22, no. 158, pp. 209--212, 1927.

\bibitem{brown2001}
L. D. Brown, T. T. Cai, and A. DasGupta, ``Interval estimation for a binomial proportion,'' \emph{Statist. Sci.}, vol. 16, no. 2, pp. 101--133, 2001.

\bibitem{mcnemar1947}
Q. McNemar, ``Note on the sampling error of the difference between correlated proportions or percentages,'' \emph{Psychometrika}, vol. 12, no. 2, pp. 153--157, 1947.

\bibitem{dietterich1998}
T. G. Dietterich, ``Approximate statistical tests for comparing supervised classification learning algorithms,'' \emph{Neural Comput.}, vol. 10, no. 7, pp. 1895--1923, 1998.

\bibitem{holm1979}
S. Holm, ``A simple sequentially rejective multiple test procedure,'' \emph{Scand. J. Statist.}, vol. 6, no. 2, pp. 65--70, 1979.

\end{thebibliography}
